% ─────────────────────────────────────────────────────────────────────
\subsection{Talent Market: A Composable Agent Ecosystems}
\label{sec:talent-market}
% ─────────────────────────────────────────────────────────────────────

The Talent Market is not merely a sourcing mechanism for agents, but a standardisation layer that enables heterogeneous agents to be integrated as interoperable components within the OMC system. While the Vessel defines \emph{how} an agent executes, the Talent Market defines \emph{how an agent becomes a compatible organizational unit}. Its primary role is to provide a uniform Talent interface through which arbitrary external agents can be transformed into plug-and-play employees governed by the OMC runtime.

\textbf{From Tools to Talents.}
% The design of the Talent Market is motivated by an analogy to tool ecosystems in single-agent systems.
In a standard LLM agent setting, external capabilities are exposed as tools or skills, each conforming to a function-level interface. The agent composes these tools to achieve complex behavior. OMC extends this paradigm from the level of individual reasoning steps to the level of organizational structure: entire agents become the composable primitives. Under this view, Talents play the same role in OMC as tools and skills do in conventional agent systems.

\begin{center}
\begin{tabular}{l|l|l}
\toprule
\textbf{Level} & \textbf{Single-Agent System} & \textbf{OMC System} \\
\midrule
Execution Unit & Tool / Skill & Talent (agent) \\
Interface & Function schema / API & Harness / Company regulations \\
Composition & Tool chaining & Organizational coordination \\
\bottomrule
\end{tabular}
\end{center}


\paragraph{Talent Interface and Packaging.}
The Talent Market defines a standardised Talent interface that specifies how an arbitrary agent can be integrated into OMC. This interface is realized through a canonical \emph{Talent template}\footnote{\url{https://github.com/1mancompany/talent-template}}, which serves as the packaging format for all Talents. A Talent package encapsulates an agent behind a uniform contract, including: (i) a role and responsibility specification, (ii) a structured input/output schema for task execution, and (iii) compatibility with the Harness communication protocol. Optional components may expose internal skills or bundle auxiliary resources such as prompts and tools.

This design separates interface from implementation, that heterogeneous agents can be wrapped into interchangeable, reusable Talents that can be seamlessly composed within the organization.

\paragraph{Three-Type Agent Supply Chain.}
The Talent Market supports three complementary sourcing channels for constructing interface-compliant Talents: (i) adapting existing open-source agent implementations, (ii) instantiating role-defined agents with assembled skills, and (iii) dynamically synthesising agents from modular skill components. These channels represent alternative construction pathways rather than distinct categories, and all produce standardised Talent packages that are indistinguishable at the interface level. We provide the detailed design and implementation of each type to Appendix~\ref{sec:talent-supply}.

\textbf{Ecosystem and Reusability.}
By enforcing a uniform interface, the Talent Market enables the emergence of a reusable and extensible agent ecosystem. Talents can be shared, recomposed, and redeployed across tasks without modification to the underlying system. As more Talents are introduced, the overall capability of OMC improves which in turn attracts more users and further expands the Talent ecosystem.

% \textbf{Evolution through Feedback.}
% The Talent Market supports continuous ecosystem evolution through feedback loops across different instantiation strategies. Skills refined during deployment can be published back to the skill marketplace, improving future Talent construction. Similarly, effective synthesized agents can be distilled into reusable personas or packaged implementations, enriching the pools used by other strategies. This bidirectional flow enables the gradual emergence of higher-quality, reusable agents independent of their original construction pathway.

% \paragraph{Hiring as Interface Binding.}
% Within this framework, hiring is reframed as a selection process over a pool of interface-compliant Talents rather than a construction pipeline. The system presents a ranked set of candidates, and the CEO selects among them based on provenance and suitability. Once selected, a Talent is bound to a Vessel and enters the organization as a standard employee, with no distinction retained regarding its origin.











% old version





% % ─────────────────────────────────────────────────────────────────────
% \subsection{Talent Market: A Three-Type Agent Supply Chain}
% \label{sec:talent-market}
% % ─────────────────────────────────────────────────────────────────────

% The Talent Market is the supply-side complement to the Vessel abstraction: while the Vessel defines \emph{how} an agent executes, the Talent Market determines \emph{what} cognitive capability is available for hire. It provides three types of agent sourcing, each representing a different method of constructing a Talent package. The three types are not ranked by quality or priority---they are parallel supply channels, and all produce standard Talent packages that enter the OMC runtime through the same pathway: loaded into a Vessel and governed by the Harness interfaces. Once hired, every agent begins at the same employee level regardless of its sourcing type.

% \paragraph{Type 1: Curated Repository Agents.}
% Type~1 agents are manually packaged from established open-source agent repositories. Each source project is selected based on community adoption metrics and benchmark performance; its core reasoning strategies, tool interfaces, and domain knowledge are distilled into the OMC Talent package format, stripping framework-specific runtime dependencies. The resulting Talents are validated against reference tasks to ensure behavioral fidelity. This sourcing method is best suited for domains where mature, battle-tested agent implementations already exist in the open-source ecosystem.

% \paragraph{Type 2: Prompt-Sourced Agents with Skill Assembly.}
% Type~2 agents start from high-quality \emph{system prompts} sourced from community-curated prompt repositories---notably the Agency-Agents collection~\cite{agencyagents}, which provides over 140 specialist personas spanning engineering, design, marketing, product, sales, project management, testing, and other divisions. Each persona defines a detailed identity, working principles, deliverable specifications, and success metrics, but carries no executable tool bindings or runtime skills. To transform a persona into a deployable Talent, OMC applies an automated \textbf{skill assembly} step: the system queries the SkillsMP marketplace~\cite{skillsmp} using the persona’s domain descriptors and capability requirements as a structured search query, retrieves compatible skills (tool bindings, behavioral principles, workflow fragments), and composes them with the source persona into a complete Talent package. This sourcing method is best suited when a well-defined role description exists but no complete agent implementation is available.

% \paragraph{Type 3: Dynamic Agent Assembly from Cloud Skills.}
% Type~3 agents are assembled entirely from modular skills retrieved from SkillsMP~\cite{skillsmp}, an open community-driven skill marketplace that aggregates agent skills from public GitHub repositories using the standardized \texttt{SKILL.md} format. Unlike Type~1 (which packages an existing implementation) or Type~2 (which starts from a curated persona), Type~3 constructs \emph{both} the persona and the skill set on demand. The HR agent provides a structured job description based on the task specification; the Talent Market then performs semantic search over skill metadata, filters for mutual compatibility and harness compliance, ranks by relevance and community score, and synthesizes a coherent agent identity from the retrieved components. This sourcing method is best suited for niche or emerging domains where neither complete agent implementations nor established persona templates yet exist, and also serves as a \emph{cold-start} mechanism for the Talent Market as a whole.

% \paragraph{Human-in-the-Loop Talent Selection.}
% Talent hiring is not fully automated: the system presents a ranked \emph{top-$k$} recommendation list to the CEO, who makes the final selection. To ensure the CEO sees a diverse range of sourcing methods, the recommendation engine composes the shortlist with an approximately 80/20 distribution---roughly 80\% of candidates from Type~1 and Type~2, and 20\% from Type~3. This ratio reflects the current relative maturity of each supply channel rather than an intrinsic quality ordering: Type~1 and Type~2 draw from repositories with established community validation, while Type~3’s dynamically assembled agents offer broader coverage but have not undergone prior community review. The CEO can inspect each candidate’s provenance---source repository for Type~1, source persona and attached skills for Type~2, constituent skills for Type~3---and approve or reject before the hiring pipeline provisions a Vessel.

% \paragraph{Cross-Type Feedback.}
% The three sourcing channels are not isolated: skills refined through deployed agents’ self-improving loops can be published back to the skill marketplace, enriching the community pool regardless of which type originally produced the agent. Similarly, a Type~3 agent whose synthesized persona proves effective across multiple projects can have that persona persisted and contributed to prompt repositories, expanding the available pool for future Type~2 sourcing.